\documentclass[10pt,a4paper]{article}

% ----------------- Paquetes -------------------%
\usepackage[T1]{fontenc}
\usepackage[utf8]{inputenc}
\usepackage[a4paper,margin=2.5cm]{geometry}
\usepackage{mathptmx}             
\usepackage{changepage} 
\usepackage{setspace}
\usepackage[spanish]{babel}
\usepackage[labelfont=bf,labelsep=space]{caption}
\usepackage{amssymb}
\usepackage{amsmath}
\usepackage{ebgaramond}
\usepackage{placeins}
\usepackage{etoolbox}
\usepackage{setspace}
\usepackage{parskip} 
\usepackage{tcolorbox}
\usepackage{needspace}
\usepackage{xparse}
\usepackage{ocgx2}
\usepackage{hyperref}
\usepackage{hypcap}
\usepackage{soul}\sethlcolor{yellow}% resaltado amarillo: \hl{texto} (Panel TCEE, ⌃H)


%%%%%%%%%%%%%%%%%%%%%%%%%%%%%%%%%%%%%%%%%%%%%%%%%%%%%%%%%%%%%%%%
% --- Fija primero tus valores globales "buenos" ---
\setstretch{1.2}                 % Interlineado global
\setlength{\parskip}{0.7\baselineskip}
\setlength{\parindent}{0pt}
\newlength{\GlobalParskip}
\setlength{\GlobalParskip}{\parskip}

\newcounter{eqnum}[subsection]
\renewcommand{\theeqnum}{\arabic{eqnum}}
\newcounter{alphaitem}
\newcounter{numitem}

%% ------------------- Recuadros----------------------%%
\newcommand{\puntosclave}[1]{%
  \begin{tcolorbox}[
    colframe=black,
    title={Puntos clave},
    before upper={\setlength{\parindent}{0.5cm}\setlength{\parskip}{0.5\baselineskip}}
  ]
  #1
  \end{tcolorbox}
}

\newcommand{\modificaciones}[1]{
  \begin{tcolorbox}[
    colback=white,
    colframe=red,
    title={Modificaciones a realizar},
    before upper={\setlength{\parindent}{0.5cm}\setlength{\parskip}{0.5\baselineskip}}
  ]
  #1
  \end{tcolorbox}
}
%% ------------------- Preguntas TEST----------------------%%
% Opciones: radio (single-choice)
\NewDocumentCommand{\OptRadio}{m m m}{%
  \noindent\CheckBox[radio,name=#1,value=#2,width=1.2ex,height=1.2ex]{}%
  \;\textbf{#2.}~#3\par
}

% Opciones: checkbox (multi-choice)
\NewDocumentCommand{\OptCheck}{m m}{%
  \noindent\CheckBox[name=#1,width=1.2ex,height=1.2ex,bordercolor={0 0 0}]{}%
  \;#2\par
}

% Pregunta single-choice (radio)
% Args: 1=id, 2=enunciado, 3=opciones, 4=solución+justificación, 5=imagen (o {}), 6=origen
\NewDocumentCommand{\PreguntaTestSingle}{m m m m m m}{%
  \par\medskip
  \noindent\textbf{#1.}~#2\par\smallskip

    % Imagen (si no es {})
    \IfBlankTF{#5}{}{%
      \begin{center}
        \includegraphics[width=0.75\linewidth]{#5}
      \end{center}
    }

    \begin{Form}
      #3
    \end{Form}

    \par\smallskip
    \switchocg{sol-#1}{\fbox{Mostrar / ocultar solución}}%
    \hfill{\footnotesize #6}\par\smallskip

    \begin{ocg}{Solución}{sol-#1}{0}
      \noindent #4
    \end{ocg}
  \par\medskip
}

% Pregunta multi-choice (checkboxes)
\NewDocumentCommand{\PreguntaTestMulti}{m m m m m m}{%
  \par\medskip
  \noindent\textbf{#1.}~#2\par\smallskip

    % Imagen (si no es {})
    \IfBlankTF{#5}{}{%
      \begin{center}
        \includegraphics[width=0.75\linewidth]{#5}
      \end{center}
    }
    \begin{Form}
      #3
    \end{Form}

    \par\smallskip
    \switchocg{sol-#1}{\fbox{Mostrar / ocultar solución}}%
    \hfill{\footnotesize #6}\par\smallskip

    \begin{ocg}{Solución}{sol-#1}{0}
      \noindent #4
    \end{ocg}
  \par\medskip
}


%% ------------------- CITA ----------------------%%
% \begin{cita}[Autor][Año][Obra] texto \end{cita}: cita sangrada y, debajo a la derecha, «AUTOR (Año), Obra». Los tres datos son opcionales.
% En el Panel TCEE: escribir \cita e Intro (el cursor va al texto; Tab pasa a Autor, Año y Obra).
\NewDocumentEnvironment{cita}{ O{} O{} O{} +b }{%
  \begin{adjustwidth}{2cm}{0pt}
  \raggedright
  #4\par
  \raggedleft
  \if\relax\detokenize{#1#2#3}\relax
    % sin firma
  \else
    #1%
    \if\relax\detokenize{#2}\relax\else\if\relax\detokenize{#1}\relax\else\space\fi(#2)\fi
    \if\relax\detokenize{#3}\relax\else\if\relax\detokenize{#1#2}\relax\else, \fi\textit{#3}\fi
  \fi
  \end{adjustwidth}
}{}



%% ------------------- Listas----------------------%%
    % --- Lista alfabética ---
    \newenvironment{la}
      {\begin{list}{\alph{alphaitem})}{%
          \usecounter{alphaitem}%
          \setlength{\leftmargin}{1cm}%
          \setlength{\parskip}{3pt}% 
          \setlength{\itemsep}{0pt}% ← espacio entre ítems
          \setlength{\parsep}{0pt}%  ← párrafos dentro de un ítem
      }}
      {\end{list}}
      
    % --- Lista numérica ---
    \newenvironment{lnum}
      {\begin{list}{\arabic{numitem}.}{%
          \usecounter{numitem}%
          \setlength{\leftmargin}{1cm}%
          \setlength{\parskip}{3pt}% 
          \setlength{\itemsep}{0pt}% ← espacio entre ítems
          \setlength{\parsep}{0pt}%  ← párrafos dentro de un ítem
      }}
      {\end{list}}
      
% ------------------- Imágenes----------------------%
    \graphicspath{{Img/}}% Rutas por defecto 
    % --- Imagen adaptativa ---
    % Registros de dimensión definidos una sola vez
    \newdimen\imgcap
    \newdimen\imgmin
    \newdimen\imgmax
    \newdimen\imgremain
    \newdimen\imgh
    \newdimen\imgneed
    
    \newcommand{\imagenfit}[3]{%
      \addvspace{10pt}% separación antes
      \par\begingroup
        % Parámetros de composición (asignaciones, no \newdimen)
        \imgcap=1\baselineskip            % alto estimado de título+fuente
        \imgmin=0.15\textheight
        \imgmax=0.15\textheight
        \imgremain=\pagegoal
        \ifdim\pagegoal=\maxdimen \imgremain=0pt\fi
        \advance\imgremain by -\pagetotal
        \advance\imgremain by -\imgcap
        % Decisión de altura destino
        \ifdim\imgremain<\imgmin
          \newpage
          \imgh=0.20\textheight
        \else
          \imgh=\imgremain
          \ifdim\imgh>\imgmax \imgh=\imgmax \fi
        \fi
        % Reservar espacio para evitar cortes del bloque completo
        \imgneed=\imgh \advance\imgneed by \imgcap
        \Needspace{\imgneed}
        % Cuerpo
        \noindent\begin{minipage}{\textwidth}
          \centering
          \captionof{figure}{\textemdash\ #2}\par\vspace{0.3em}% título arriba
          \includegraphics[width=\linewidth,height=\imgh,keepaspectratio]{\detokenize{#1}}\par
          \vspace{0.3em}\small\textit{Fuente: }#3% fuente abajo
        \end{minipage}%
      \endgroup
      \par\addvspace{10pt}%
    }

%------------ Bloque de RECUADRO ESPECIAL -------------%
    \newenvironment{specialblock}[1]{%
      \begin{quote} % crea sangría a izquierda y derecha
      \small % texto más pequeño
      \noindent\textbf{#1}\\[-0.3em] % título en negrita
      \rule{\linewidth}{0.4pt}\\[0.6em]}{
      \end{quote}}
      
%-------------------------- Portada -----------------------%
\newcommand{\oldtitle}[1]{\def\OldTitle{#1}}
\newcommand{\changerationale}[1]{\def\ChangeWhy{#1}}

\newcommand{\changebox}{%
\begin{tcolorbox}[title={Historial de cambios del tema}]
\textbf{Título anterior:} \OldTitle\par
\textbf{Motivo del cambio:} \ChangeWhy
\end{tcolorbox}
}

%-------------------------- Author Font -----------------------%
    \newcommand{\authorfont}[1]{{\fontfamily{EBGaramond-TLF}\selectfont #1}}

%-------------------------- Anexos -----------------------%
    \makeatletter
    \newcounter{anexo}
    \renewcommand{\theanexo}{\Roman{anexo}}
    \newcommand{\anexo}[1]{%
      \clearpage
      \phantomsection
      \refstepcounter{anexo}%
      \noindent\textbf{Anexo \theanexo.- }#1\par\medskip
      \addcontentsline{stc}{section}{Anexo \theanexo.- #1}%
    }
    \makeatother

    %-------------------------- Lista de figuras (personal) -----------------------%
\makeatletter
\newcommand{\MyListOfFigures}{%
  \clearpage
  \phantomsection
  {\centering\bfseries\Large Índice de figuras\par}%
  \medskip
  % Entrada en nuestro índice de contenido (stc)
  \addcontentsline{stc}{section}{Índice de figuras}%
  % Recuperación del listado (sin cabecera propia)
  \begingroup
    \parindent=0pt\parskip=2pt
    \@starttoc{lof}%
  \endgroup
}
\makeatother

%-------------------------- Bibliografía -----------------------%
\makeatletter
% Contador para las referencias
\newcounter{myref}
% Cabecera de la bibliografía, entra en el índice 'stc'
\newcommand{\MyBibliography}{%
  \clearpage
  \phantomsection
  {\centering\bfseries\Large Bibliografía y referencias\par}%
  \medskip
  \addcontentsline{stc}{section}{Bibliografía y referencias}%
  \setcounter{myref}{0}%
}
\newcommand{\MyBibItem}[4]{%
  \refstepcounter{myref}%
  \noindent[\themyref]\ #1.\ \emph{#2}. #3. #4\par
}
\makeatother

%--------- Establecer cuatro niveles de índice --------%
    \setcounter{tocdepth}{4}   
    \setcounter{secnumdepth}{4}% numera hasta \paragraph

%%%%%%%%%%%%%%%%%%%%%%%%%%%%%%%%

\makeatletter

    %--------- Interlineado Global --------%
    \edef\GlobalBaseLineStretch{\baselinestretch}
    
    %--------- Espaciado de seccinones --------%
    \renewcommand\section{\@startsection{section}
      {1}{\z@}
      {2.5ex \@plus 1ex \@minus .2ex} % espacio antes
      {1.5ex \@plus .2ex}             % espacio después
      {\normalfont\Large\bfseries}    % formato del título
    }
    \renewcommand\subsection{\@startsection{subsection}{2}{\z@}
      {3ex \@plus 0.8ex \@minus .2ex}
      {1.5ex \@plus .2ex}
      {\normalfont\large\bfseries}
    }
    \renewcommand\subsubsection{\@startsection{subsubsection}{3}{\z@}
      {3ex \@plus 0.5ex \@minus .2ex}
      {1ex \@plus .2ex}
      {\normalfont\normalsize\bfseries}
    }
    \renewcommand\paragraph{\@startsection{paragraph}{4}{\z@}%
    {-2ex \@plus -0.5ex \@minus -.2ex}% espacio antes
    {0.8ex \@plus .2ex}% espacio después
    {\normalfont\normalsize\itshape}}% ← cursiva en lugar de negrita

    %--------- Ecuaciones --------%   
    % Variables globales para almacenar títulos y notas
    \gdef\leq@current@section{}%
    \gdef\leq@current@subsection{}%
    \gdef\leq@last@printed@section{}%
    \gdef\leq@last@printed@subsection{}%
    
    % Comando para añadir nota (invisible en el cuerpo)
    \newcommand{\eqnote}[1]{%
      \addcontentsline{leq}{leqnote}{#1}%
    }
    
    % Comando para ecuaciones
    \newcommand{\eqblock}[2]{%
      % Verificar si necesitamos imprimir el título de sección
      \ifx\leq@current@section\leq@last@printed@section\else
        \addcontentsline{leq}{leqsection}{\leq@current@section}%
        \global\let\leq@last@printed@section\leq@current@section
        \gdef\leq@last@printed@subsection{}% Reset subsección al cambiar sección
      \fi
      % Verificar si necesitamos imprimir el título de subsección
      \ifx\leq@current@subsection\leq@last@printed@subsection\else
        \ifx\leq@current@subsection\@empty\else
          \addcontentsline{leq}{leqsubsection}{\leq@current@subsection}%
          \global\let\leq@last@printed@subsection\leq@current@subsection
        \fi
      \fi
      % Ecuación
      \refstepcounter{eqnum}%
      \begin{equation*}
        #1\tag{E.\theeqnum}
      \end{equation*}%
      \addcontentsline{leq}{leqt}{[E.\theeqnum] -- #2}%
      \addcontentsline{leq}{leqm}{#1}%
    }
    
    %%% ----- Índice de ecuaciones ----------%%%
    % Tamaño para []
    \newcommand{\leqIndexFont}{\normalsize}
    \newcommand{\leqTitleFont}{\tiny}
    \newcommand{\leqNoteFont}{\tiny}
    % Cabecera de sección 
    \newcommand*\l@leqsection[2]{%
      \addvspace{25pt}% Mayor separación antes de nueva sección
      \noindent\leqIndexFont
      \makebox[\linewidth]{\rule{.38\linewidth}{0.4pt}\ \ \textbf{  \MakeUppercase{#1}}\ \ \rule{.38\linewidth}{0.4pt}}%
      \par\vspace{-10pt}
    }
    % Cabecera de subsección 
    \newcommand*\l@leqsubsection[2]{%
      \par\addvspace{15pt}%
      \noindent\leqIndexFont #1
      \par\vspace{-7pt}
    }
    % Título de ecuación 
    \newcommand*\l@leqt[2]{%
    \par\addvspace{10pt}%
      \par\noindent\leqTitleFont #1\par
    }
    % Miniatura de ecuación
    \newcommand*\l@leqm[2]{%
      \vspace{0.3\baselineskip}%
      \par\scriptsize\noindent\hspace*{1.5em}\(#1\)\par%
    }
    % Nota de ecuación 
    \newcommand*\l@leqnote[2]{%
      \vspace{0.2\baselineskip}%
      \par\noindent\leqNoteFont\hspace*{1.5em}\textbf{Nota.--} #1\par\vspace{0.5\baselineskip}%
    }
    % Generación del índice
    \newcommand{\mylistofequations}{%
      \clearpage
      \phantomsection
      % Título visual de la lista (ajústalo a tu gusto):
      {\centering\bfseries\Large Índice de ecuaciones\par}
      % Entrada en nuestro índice 'stc':
      \addcontentsline{stc}{section}{Índice de ecuaciones}%
      % Llamada a tu lista real (usa el nombre exacto de tu macro):
      \@starttoc{leq}%
    }
    
    %--------- Captura de títulos de sección --------%
    \let\leq@original@section\section
    \let\leq@original@subsection\subsection
    % Flag para saber si estamos en una sección asterisco
    \newif\ifLeqStarSection
    \LeqStarSectionfalse
    
    % Redefinir \section CON CONTROL DE ASTERISCO
    \renewcommand{\section}{%
      \@ifstar{\leq@section@star}{\@dblarg\leq@section@nostar}%
    }
    \def\leq@section@star#1{%
      \global\LeqStarSectiontrue
      \gdef\leq@current@section{#1}%
      \gdef\leq@current@subsection{}%
      \leq@original@section*{#1}%
      \global\LeqStarSectionfalse
    }
    \def\leq@section@nostar[#1]#2{%
      \global\LeqStarSectionfalse
      \gdef\leq@current@section{#2}%
      \gdef\leq@current@subsection{}%
      \leq@original@section[#1]{#2}%
    }
    % Redefinir \subsection
    \renewcommand{\subsection}{%
      \@ifstar{\leq@subsection@star}{\@dblarg\leq@subsection@nostar}%
    }
    \def\leq@subsection@star#1{%
      \global\LeqStarSectiontrue
      \gdef\leq@current@subsection{#1}%
      \leq@original@subsection*{#1}%
      \global\LeqStarSectionfalse
    }
    \def\leq@subsection@nostar[#1]#2{%
      \global\LeqStarSectionfalse
      \gdef\leq@current@subsection{#2}%
      \leq@original@subsection[#1]{#2}%
    }

    % ----- Restauración de valores Globales de estilo ----------%
    % Flags
    \newif\ifInFootnote
    \InFootnotefalse
    \pretocmd{\@footnotetext}{\InFootnotetrue}{}{}
    \apptocmd{\@footnotetext}{\InFootnotefalse}{}{}
    % Profundidad de listas (soporta anidadas)
    \newcount\MyListDepth \MyListDepth=0
    \AtBeginEnvironment{la}{\advance\MyListDepth by 1}
    \AfterEndEnvironment{la}{\advance\MyListDepth by -1}
    \AtBeginEnvironment{lnum}{\advance\MyListDepth by 1}
    \AfterEndEnvironment{lnum}{\advance\MyListDepth by -1}
    \AtBeginEnvironment{itemize}{\advance\MyListDepth by 1}
    \AfterEndEnvironment{itemize}{\advance\MyListDepth by -1}
    \AtBeginEnvironment{enumerate}{\advance\MyListDepth by 1}
    \AfterEndEnvironment{enumerate}{\advance\MyListDepth by -1}
    % Restauración diferida: hazlo al INICIO del próximo párrafo real
    \newif\if@pendingRestore
    \@pendingRestorefalse
    \newcommand{\RequestRestore}{\global\@pendingRestoretrue}
    \everypar{%
      \if@pendingRestore
        \global\@pendingRestorefalse
        \ifInFootnote\else
          \ifnum\MyListDepth=0
            \setlength{\parskip}{\GlobalParskip}%
            \setstretch{\GlobalBaseLineStretch}%
          \fi
        \fi
      \fi
    }
    % Al final de bloques:
    \AfterEndEnvironment{equation}{\RequestRestore}
    \AfterEndEnvironment{equation*}{\RequestRestore}
    \AfterEndEnvironment{la}{\RequestRestore}
    \AfterEndEnvironment{lnum}{\RequestRestore}
    % Al final de macros:
    \apptocmd{\eqblock}{\RequestRestore}{}{}
    \apptocmd{\imagenfit}{\RequestRestore}{}{}
    
\makeatother

% ===================== ToC propio con 4 niveles (único bloque) =====================
\makeatletter

% --- Escritores: añadimos una línea al .stc en cada cabecera numerada ---
% Guardamos las versiones actuales para llamar al comportamiento normal
\let\MyTOC@orig@section\section
\let\MyTOC@orig@subsection\subsection
\let\MyTOC@orig@subsubsection\subsubsection
\let\MyTOC@orig@paragraph\paragraph

% SECTION
\renewcommand{\section}{%
  \@ifstar{\MyTOC@section@star}{\@dblarg\MyTOC@section@nostar}%
}
\def\MyTOC@section@star#1{%
  \MyTOC@orig@section*{#1}% sin entrada al .stc
}
\def\MyTOC@section@nostar[#1]#2{%
  \MyTOC@orig@section[{#1}]{#2}% comportamiento normal (escribe al .toc)
  % Escribimos también nuestra línea al .stc (texto con \numberline)
  \addcontentsline{stc}{section}{\protect\numberline{\thesection}#2}%
}

% SUBSECTION
\renewcommand{\subsection}{%
  \@ifstar{\MyTOC@subsection@star}{\@dblarg\MyTOC@subsection@nostar}%
}
\def\MyTOC@subsection@star#1{%
  \MyTOC@orig@subsection*{#1}%
}
\def\MyTOC@subsection@nostar[#1]#2{%
  \MyTOC@orig@subsection[{#1}]{#2}%
  \addcontentsline{stc}{subsection}{\protect\numberline{\thesubsection}#2}%
}

% SUBSUBSECTION
\renewcommand{\subsubsection}{%
  \@ifstar{\MyTOC@subsubsection@star}{\@dblarg\MyTOC@subsubsection@nostar}%
}
\def\MyTOC@subsubsection@star#1{%
  \MyTOC@orig@subsubsection*{#1}%
}
\def\MyTOC@subsubsection@nostar[#1]#2{%
  \MyTOC@orig@subsubsection[{#1}]{#2}%
  \addcontentsline{stc}{subsubsection}{\protect\numberline{\thesubsubsection}#2}%
}

% PARAGRAPH
\renewcommand{\paragraph}{%
  \@ifstar{\MyTOC@paragraph@star}{\@dblarg\MyTOC@paragraph@nostar}%
}
\def\MyTOC@paragraph@star#1{%
  \MyTOC@orig@paragraph*{#1}%
}
\def\MyTOC@paragraph@nostar[#1]#2{%
  \MyTOC@orig@paragraph[{#1}]{#2}%
  % Si no numeras los \paragraph, cambia \theparagraph por texto plano sin \numberline
  \addcontentsline{stc}{paragraph}{\protect\numberline{\theparagraph}#2}%
}

% --- Impresor del índice propio (lee el .stc) ---
\newcommand{\MyTOC}{%
  \par\bigskip
  {\centering\bfseries\Large Índice de contenido\par}%
  \medskip
  \begingroup
    \parindent=0pt\parskip=2pt
    \@starttoc{stc}%
  \endgroup
}

\makeatother
% ===================== fin ToC propio =====================


%%%%%%%%%%%%%%


% --- Datos del tema ---
\title{Tema 3.A.36 -- Política Monetaria (I)\\
\large El diseño y la instrumentación de la Política Monetaria}
\author{Marcelo Ortega}
\date{Enero 2026}
\newcommand{\TemaCode}{3A36}

\oldtitle{Tema 3.A.36. Política Monetaria (I). El diseño y la instrumentación de la Política Monetaria}
\changerationale{Sin cambios.}

\begin{document}


\maketitle
\changebox

\newpage

% --- Página de resumen y modificaciones ---
\puntosclave{ %% Escribir puntos clave Apuntes CECO
     
    }
\vspace{1cm}

\modificaciones{
    Hay que ver cómo podemos introducir la idea (si se puede) de qué objetivos como Price-Level-Targeting y Nominal-GDP-Targeting se conciben como políticas regladas sobre variables instrumentales, instrumentos o objetivos intermedios que persiguen dichos objetivos finales. A diferencia del Inflation Targeting cuyo objetivo=regla que se sigue.

    
    }

\newpage
\MyTOC
\newpage

\mylistofequations
\clearpage

\MyListOfFigures
\clearpage


\pagenumbering{arabic}
\setcounter{page}{1}


\section{Introducción}



\subsection{Contextualización}


\subsubsection{Relevancia}


\subsubsection{Problemática}



\subsection{Estructura de la exposición}



\clearpage
\section{El Dinero: Creación de la Oferta Monetaria ($M^s$)}
La exposición del diseño e intrumentos de los que dispone la Autoridad Monetaria para llevar a cabo su objetivo, lo que se conoce como Política Monetaria, exige que introduzcamos brevemente el proceso de creación de la Oferte Monetario, es decir, de los agregados monetarios o el dinero en sentido amplio. Desde una perspectiva económica, si en el [\authorfont{Ver Tema 3.A.35}] exponíamos la Teoría de la Demanda de Dinero, es decir, el interés o deseo que mueve a los agentes para demandar dinero, toca ahora analizar la otra cara de la moneda, la oferta.

\subsection{Demanda de dinero: Breve recapitulación de la \textit{utilidad} o funciones del dinero}
Como sabemos, el dinero no tiene una definición única y completa, por lo que, brevemete, conviene enunciar las razones que determinan la demanda de dinero por parte de los agentes que son, de base, sus principales cualidades:
\begin{lnum}
    \item \textbf{Medio de pago}. El dinero se define como el método generalmente aceptado para la realización de transacciones y la cancelación de deudas (en palabras de CLOWER, uno de los pioneros de la Teoría Monetaria: \textit{el dinero compra bienes y los bienes compran dinero, pero los bienes no compran bienes}). Bajo este pretexto se entiende que los agentes demandan dinero porque les proporciona utilidad de forma indirecta por su capacidad de comprar bienes (que les proporcionan utilidad de manera directa);
    \item \textbf{Motivo precaucióno depósito de valor}. Los agentes demandan dinero, en condiciones normales, por se un bien que preserva su valor en el tiempo y puede ser empleado para realizar transacciones en el futuro. Por tanto, vemos como la confianza y las expectativas de los agentes juegan un papel fundamental en la demanda de dinero; siendo de hecho el determinante más destacado en la Teoría Monetaria Moderna; y,
    \item \textbf{Motivo especulación}. Los agentes demandan bonos (activos nominales alternativos) o dinero en función del tipo de interés nominal ($i_t$), buscando percibir la mayor rentabilidad posible.
\end{lnum}

Además, el dinero tiene la utilidad de servir de \textbf{unidad de cuenta estandarizada}, pues actúa de facilitador de medida del precio de los bienes y los servicios. 

\subsection{El proceso de creación del dinero: la Oferta Monetaria}
Dicho esto, adentremonos en la primera parte del Tema: la creación de la Oferta Monetaria. Para comprender el diseño de objetivos y los instrumentos de los que dispone la PM, responderemos primero a las siguientes cuestiones: ¿qué es el dinero en la actualidad? ¿qué son los agregados monetarios? ¿qué incidencia tiene la Autoridad Monetaria en la oferta? ¿cómo se crea el dinero?

\subsubsection{Agregados Monetarios}
Lo primero que conviene realizar es una definición operativa que nos sirva para delimitar los activos que forman parte del dinero, es decir, como este este concepto se proyecta en la realidad económica. Esto es lo que se conoce como los diferentes \textbf{agregados monetarios} y son el primer paso para entender los componentes de la oferta. 

Estos agregados han sufrido una transformación constante y sustancial a lo largo de la historia del dinero (en paralelo a su propia definición y conceptualización, actualmente: \textit{fiat money}) que es lo que, comúnmente se concoe como \textbf{innovaciones financieras}. Estas innovaciones exigen su traslación teórica para obtener una mayor comprensión de los fenómenos que se observan en el día a día. En la actualidad, la teoría económica propone la siguiente clasificación de los agregados monetarios, cuya diferencia principal podríamos decir que se reduce al grado de liquidez/disponibilidad en la que se presenta el dinero:
\begin{lnum}
    \item \textbf{Base Monetaria ($M_0$)}. Equivalente a los billetes y monedas en circulación, y las reservas bancarias (BC y Comerciales). Es típicamente el agregado de menor cuantía, es decir el más estrecho, y que proyecta más adecuadamente la labor del Banco Central como monopolista en la creación de dinero y por tanto en su control, motivo por el cual este agregado ha jugado históricamente un papel importante en la conducción de la política monetaria;
    \item \textbf{Mercado Monetario}. Los siguientes agregados son los que conforman el mercado monetario en su sentido actual.
    \begin{lnum}
        \item \textbf{Aregado Monetario estrecho ($M_1$)}. Se trata del efectivo en circulación (posesión privada) y la suma de depósitos a la vista;\footnote{%
            Los depósitos a la vista son productos financieros (como cuentas corrientes o cuentas de ahorro) que te permiten guardar dinero en un banco y disponer de él en cualquier momento sin previo aviso, ofreciendo gran liquidez para gastos diarios y permitiendo realizar operaciones como transferencias, domiciliaciones y el uso de tarjetas de débito. Aunque ofrecen una rentabilidad (intereses) por tener el dinero parado, esta suele ser baja debido a su disponibilidad inmediata
        }
        \item \textbf{Agregado monetario intermedio ($M_2$)}. Se constituye por la suma del dinero en sentido estricto ($M_1$) y los depósitos de ahorro;\footnote{
        Los depósitos de ahorro, que no cuentas corrientes, están conformados por los siguientes tipos (estríctamente contenidos en $M_2$):
        \begin{lnum}
            \item \textbf{Depósitos a plazo de hasta dos años}. Depósito contratado por un plazo de tiempo de hasta dos años. Si el contrato contiene una cláusula que así lo indique, el importe del depósito solo se puede retirar antes de que expire el plazo con autorización de la entidad financiera y con una penalización en su remuneración; y,
            \item \textbf{Depósitos disponibles con preaviso}. Depósitos no susceptibles de transformación en medios de pago, sin plazo de vencimiento acordado, que no pueden convertirse en efectivo sin un período de preaviso, antes del cual la conversión en efectivo no es posible, o bien lo es únicamente con una penalización o restricción significativa.
        \end{lnum}
        }
        \item \textbf{Agregado monetario en sentido amplio ($M_3$)}. Comprende los pasivos incluidos en $M_2$ más los depósitos a plazo mayor a dos años y otros pasivos del balance de Bancos o Instituciones Financieras;\footnote{%
            Entre otras, las cesiones temporales, las participaciones en fondos del mercado monetario e instrumentos del mercado monetario y los valores de renta fija de hasta dos años, emitidos por las instituciones financieras monetarias.
        }
        \item \textbf{Agregado monetario total ($M_4$)}. Se trata del agregado monetario menos líquido, ya que incluye el $M_3$ más las emisiones de bonos, letras y pagarés y los activos emitidos por las instituciones financieras que puedan ser considerados cuasi dinero. No obstante, no está aceptado de forma generalizada ya que, entre otros, el FMI no lo utiliza.
    \end{lnum}
\end{lnum}

De lo anterior se puede deducir fácilmente que el dinero es un concepto muy amplio que se expande a partir de su forma más simple y tradicional: el \textbf{efectivo en manos del público}.

\subsubsection{Proceso de creación de los agregados monetarios: la Oferta Monetaria}
Además, aunque ya hemos delimitado ligeramente los agregados en los que incide directamente el Banco Central y aquellos que forman parte del sistem de creación \textit{endógena} de la Oferta Monetaria, cabe repasar cuál es el proceso de creación de cada uno de estos agregados brevemente. 

\paragraph{Base Monetaria: Banco Central}
El primer elemento del proceso recae en manos del emisor del dinero legal, aquel que actúa en régimen de Monopolio: el Banco Central. Éste incide directamente sobre la Base Monetaria que, como hemos visto está formada por el efectivo en manos del público ($E$) y las reservas totales ($TR$). Las reservas totales pueden dividirse en dos tipos: (i) las Reservas Legales ($rr$), es decir, el porcentaje sobre los depósitos que los Bancos deben mantener legalmente; y, (ii) el exceso de reservas ($er$, muy importante en la actualidad por ser el coeficiente de reservas muy bajo).

Analíticamente esto puede expresarse como:

\eqblock{
\begin{aligned}
    \text{Base Monetaria: }(M_0)\quad \overbrace{RR+ER}^{\text{\scriptsize RT de Bancos}}+C\quad\Longrightarrow\quad  \boxed{M_0=(rr+ex+c)D}
\end{aligned}
}{Base Monetaria. Proceso de creación de la Oferta Monetaria}

Desde una perspectiva contable, la Base Monetaria constituye el pasivo del Banco Central y se encuentra respaldado por sus activos. Por lo tanto, tenemos como el Balance del Banco Central se puede expresar como:

\imagenfit{Balance_BC.png}{Balance simplificado del Banco Central}{Elaboración propia}

\paragraph{Multiplicación de la BM: Mercado financiero}
Ahora bien, ¿cómo se proyecta y expande la Base Monetaria sobre el sistema monetaria? A través de los intermediarios financieros; por eso, como veremos estos juegan un papel fundamental en la transmisicón einstrumentalización de la Política Monetaria.

La dinámica es sencilla: los Bancos tienen la capacidad de captar depósitos y de crear crédito, estando en última instancia limitados por las restricciones legales impuestas por el coeficiente mínimo de reservas. Cualquier que sea la dirección que se analice: crean depósitos y conceden crédito o conceden crédito y crean depósitos (en un mismo Banco o en diferentes), el resultado es el mismo: a partir de la fracción que no mantienen en reservas, se genera una expansión de los depósitos que, en esencia, constituye una expansión de la oferta monetaria (i.e. creación de dinero por la expansión de los agregados monetarios). 

Desde una perspectiva contable, el Balance de un intermediario financiero (Banco Comercial) puede simplificarse como sigue:

\imagenfit{Balanca_BCC.png}{Balance simplificado de un Banco Comercial}{Elaboración propia}

La concesión del crédito queda supeditado por la siguiente ecuación: $CC=(rr+er)D$. Además, cabe destacar una cosa fundamental: el dinero creado a través del sistema bancario es igualmente destruible. Es decir, su razón de ser es una cuestión meramente contable por lo que el deterioro contable de los créditos trae consigo el deterioro en la misma cuantía de los depósitos generados: el dinero se crea y se destruye contablemente. De la misma forma que el dinero puede ser creado sin la necesidad de intervención del banco central, el dinero puede ser destruido de forma completamente independiente a la autoridad monetaria. Cada vez que un hogar paga a un banco comercial el préstamo que le fue concedido, el dinero que había sido previamente creado es destruido.

\paragraph{Sistema Monetario: Balance}
Por lo tanto, como vemos, la oferta monetaria constituye el pasivo del Sistema Monetario (BC y Bancos Comerciales) y este está constituido por el efectivo en manos del público ($E$) y los depósitos ($D$). Estos pasivos están respaldados a su vez por los activos del Sistema Monetario, tanto los que están en posesión del Banco Central como los que están en posesión de los Bancos Comerciales. Por tanto, podemos presentar un balance simplificado general del Sistema Monetario como:

\imagenfit{Balance_SM.png}{Balance (simplificado) del Sistema Monetario. Visión global}{Elaboración propia}

\textbf{Nota.-} EL tamaño o proporción de los diferentes componentes no guardan relación alguna con el sistema real. Se hace una aproximación simplificada con el fin de ilustrar la capacidad de creación/expansión de oferta monetaria que tiene el propio Sistema Monetario. Es lo que se conoce como endogeneización de la creación de dinero. Un proceso detallado de cómo este proceso se realiza e interactúan los elementos entre sí se puede ver en el [\authorfont{Ver Anexo \ref{anx:sistemamonetario}}], así como el Balance detallado del Banco Central.

\paragraph{Teoría del Multiplicador Monetario}
Como vemos, de lo aterior se deduce que el Banco Central no controla directamente la Oferta Monetaria; sino que controla la Base Monetaria directamente y son los Bancos Comerciales (o intermediarios financieros) los que expanden la oferta monetaria a través de la expansión de los depósitos.

La capacidad de creación/expansión de oferta monetaria que ostentan los Bancos Comerciales o, desde la perspectiva del \textit{policy-maker}, la capacidad del Banco Central para determinar la oferta de dinero agregada a través del control de la base monetaria se conoce como la Teoría del Multiplicador Monetario y puede cuantificarse a través de la siguiente expresión:

\eqblock{
    \begin{aligned}
        \boxed{mm=\frac{1+c}{\underbrace{rr+er}_{r}+c}}\quad\equiv\quad \frac{\delta M_1}{\delta M_0}
    \end{aligned}
}{Multiplicador monetario. Proceso de creación de la Oferta Monetaria}

Por ejemplo, a partir de la inyección inicial del Banco Central (de dinero en efectivo o reducción del coeficiente de reservas), las interacciones entre la Base Monetaria y la Oferta Monetaria ($M_1$) instrumentada a través de los coeficientes de efectivo y reservas totales, se puede ver como se expande en mayor cuantía la Oferta Monetaria presente en el Sistema. 

Ahora bien, este proceso sencillo enfrenta enormes dificultades cuando se traslada a la realidad. En especial, para que esto se cumpliese, el coeficiente de reservas debería de ser una restricción vinculante para todos los bancos comerciales, hecho que rara vez se cumple, ya que el coeficiente de caja obligatorio que aplican en la actualidad los bancos centrales es muy bajo, por lo que los bancos comerciales suelen decidir tener un volumen de reservas superior al que dicta dicho coeficiente. Además, una expansión de los agregados monetarios (más cantidad de efectivo) no tiene porque traducirse en una expansión mayor de la Oferta Monetaria (por ejemplo si nos encontramos en la \textit{trampa de liquidez}).

El proceso de creación de dinero recién descrito está limitado por diferentes fuerzas, cuyo determinante en última instancia es la autoridad monetaria. En la práctica, el proceso de creación del dinero responde a una lógica de oferta y demanda en el mercado, en el que los demandantes son los agentes económicos y los oferentes son las entidades de crédito. Así, los determinantes de la Oferta Monetaria pueden resumirse en:
\begin{lnum}
    \item \textbf{Límites a la cantidad de dinero que los bancos pueden prestar}. En un mercado competitivo, si un banco desea aumentar los préstamos que concede, deberá de reducir el tipo de interés que pide a cambio de los mismos. Debido a que los bancos buscan obtener cierta rentabilidad, no pueden incrementar el volumen de préstamos sin limites. Por otra parte, el aumento del volumen de préstamos incrementa la exposición al riesgo del banco, lo que puede poner en pel igro la continuidad de la entidad. Esta razón también limita la creación de dinero;
    \item \textbf{Comportamiento de hogares y empresas}. La creación de dinero también depende de los hogares y empresas que reciben los depósitos creados por los bancos comerciales. Vamos a estudiar dos escenarios, cada uno de ellos con implicaciones muy distintas sobre la creación de dinero, la actividad económica y las presiones inflacionarias de la economía:
    \begin{la}
        \item En este primer escenario, el nuevo depósito creado aumenta las tenencias de saldos monetarios reales de los agentes. Éstos intentarán ajustar este exceso de tenencias a través de la adqu isición de nuevos bienes y servicios. Ello incrementa la demanda agregada de la economía e impulsa las presiones inflacionistas;
        \item En un segundo escenario, el dinero que ha sido recientemente creado con el depósito es directamente destruido al ser utilizado para pagar previas deudas que el hogar o la empresa mantenía. Por tanto, el aumento transitorio en la oferta de dinero no llega a tener efectos reales sobre la economía.
    \end{la}
    \item \textbf{Política Monetaria aplicada por el Banco Central}. El tipo de interés fijado por el banco central influye sobre los tipos de interés demandados por los bancos comerciales como contraprestación a los préstamos que realizan. En consecuencia, políticas monetarias de carácter expansivo generarán un incremento de la oferta monetaria en la economía a través del aumento en el volumen de préstamos que concederán los bancos comerciales como consecuencia de la caída de los tipos de interés. Este incremento en la oferta monetaria se traducirá en un aumento de la demanda agregada, lo que impulsará las presiones inflacionistas y el nivel de inflación esperado en el futuro. Otro elemento de la política monetaria que ahora mismo no es determinante, pero podría llegar a serlo en determinadas circunstancias es el coeficiente de reserva, al potencialmente fijar la cantidad mínima de dinero en efectivo que los bancos han de mantener como porcentaje de sus pasivos
\end{lnum}

En conclusión, a la hora de diseñar la Política Monetaria debemos tener muy presente que el Banco Central controlará más o menos el funcionamiento del Sistema Monetario en la medida que la demanda de dinero sea estable (cumpliéndose la relación establecida por el multiplicador monetario).

\clearpage
\section{Diseño de la Política Monetario}
\puntosclave{
    La política monetaria puede perseguir distintos objetivos. Tanto la teoría como la evidencia empírica indican que su mayor contribución a la estabilidad económica es a través de la estabilidad de precios. 
    
    La discrecionalidad puede conducir a resultados subóptimos como consecuencia de los problemas de inconsistencia dinámica. Por ello, la adopción de reglas o de un conjunto de principios sistemáticos de conducta otorgan al banco central una credibilidad que le permite anclar las expectativas de inflación. 
    
    La mayoría de bancos centrales de las economías desarrolladas siguen un régimen del tipo \textit{inflation targeting}. La regla de conducción de la política monetaria bajo este régimen que más popularidad ha alcanzado es la regla de TAYLOR.
    
    En la actualidad, la política monetaria se enfrenta a importantes retos, como es el \textit{Effective Lower Bound}, el cual limita su capacidad para mantener una inflación cercana al objetivo (por anclarse a la baja) cuando el tipo de interés ya ha alcanzado su límite nominal inferior.
}

Ahora que conocemos cúal es el proceso de creación de la oferta monetaria, pasemos a analizar en profundidad el papel que desempeña la Atoridad Monetaria en el.

En particular, ¿qué es la Política Monetaria y qué forma toma? La política monetaria se define como la actuación de las autoridades monetarias encaminada a controlar la cantidad y/o el precio del dinero con el objetivo de atender objetivos de carácter macroeconómico que van desde la estabilidad de precios hasta la estabilidad cambiaría, pasando por la consecución de un detenninado nivel de actividad económica medido a partir de indicadores macroeconómicos como el PIB o la tasa de desempleo.

Entonces, ¿cómo alcanza el Banco Central los diferentes objetivos que puede tener la Política Monetaria? Tal y como introducimos en la Teoría de la Demanda de Dinero, uno de los factores determinantes que condicionan la incidencia de las variables nominales sobre las reales es la existencia de rigideces (en este caso, nominales).\footnote{
Por ejemplo, variaciones en el tipo de interés nominal (variable controlada por la autoridad monetaria) que no son inmediatamente acompañadas por variaciones proporcionales en la inflación generan cambios en el tipo de interés real, el cual afecta a las decisiones de consumo, inversión y demanda de saldos monetarios reales por parte de los agentes, entre otros. 
} Esto implica que las autoridades monetarias pueden, en principio, incidir sobre el ciclo en el nivel de variables macroeconómicas reales como el nivel de desempleo o del PIB.

No obstante, tal y como establece la Regla de TINBERGEN (1952), es imposible alcanzar dos objetivos independientes de política económica con un único instrumento. Necesitamos, por tanto, especificar \textbf{qué objetivo queremos alcanzar}: la estabilidad real medida a través de cierta variable macroeconómica o la estabilidad nominal. 

Haciendo uso de indicadores, reglas y procedimientos preestablecidos, la Autoridad Monetaria (de ahora en adelante, AM) podrá decidir cómo y cuándo utilizar los instrumentos de los que dispone para alcanzar sus objetivos. Pero, ¿qué objetivo queremos alcanzar? 

Obviando las experiencias relativas a la persecución de objetivos macroeconómicos reales como el pleno empleo o el crecimiento sobre el PIB potencial (muy frecuente durante los años 60's y 70's, lo que se conoce como \textit{políticas activistas} y que puede ser consultado en [\authorfont{Ver Tema 3.A.35 ICEX-CECO}]) debido al rechazo generalizado que la efectividad de la Política Monetaria tiene sobre estas en el largo plazo; podemos considerar que el objetivo último perseguido por la Autoridad Monetaria es, por ejemplo, la estabilidad de precios.

Entonces, ¿cómo diseñar una política que consiga alcanzar este objetivo?

\subsection{Diseño tradicional de la Política Monetaria}
Presentamos en primer lugar el diseño tradicional de la Política Monetaria, consistente en la persecución de un objetivo intermedio instrumental para alcanzar nuestro objetivo final. Se podría representar mediante el siguiente esquema operativo:

\imagenfit{Diseño_Tradicional_PM.png}{Flujo operativo en la ejecución de la Política Monetaria}{Tema 3.A.34. Apuntes Víctor García Begue}

La autoridad monetaria dispone de unos instrumentos que se representan mediante una serie de variables (operativas), estas variables operativas inciden directamente en el comportamiento de otras variables (de entre los cuales fijamos un objetivo intermedio) que a su vez mantienen relaciones estructurales (o creemos) con otras variables, de entre las cuales se encuentra nuestra variable objetivo.

\subsubsection{Segunda etapa: objetivo final}
Supongamos que esta autoridad persigue la estabilidad de precios. Esto podría representarse a través de la ecuación cuantitativ del dinero cómo:

\eqblock{
\begin{aligned}
    \dot m_t+\dot v_t=\underbrace{\dot p_t+\dot y_t}_{\text{\scriptsize Objetivo último}}
\end{aligned}
}{Objetivo final: estabilidad de precios. Diseño tradicional de la Política Monetaria}

Ahora bien, los instrumentos de los que disponemos no son infinitos, de hecho, son relativamente escasos. Por tanto, la clave radica en encontrar un objetivo intermedio que esté relacionado con nuestro objetivo final y sobre el que podamos incidir directamente con los instrumentos de los que disponemos: nuestra variable operativa.

\subsubsection{Primera etapa: variable operativa}
¿Sobre qué podemos incidir con nuestros instrumentos? Si nuestro objetivo es la estabilidad de precios y consideramos valida la ecuación cuantitativa presentada, nuestro objetivo intermedio resulta claro: los agregados monetarios ($\dot m_t$).

\eqblock{
\begin{aligned}
    \underbrace{\dot m_t}_{\text{\scriptsize Objetivo intermedio}}+\dot v_t=\underbrace{\dot p_t+\dot y_t}_{\text{\scriptsize Objetivo último}}
\end{aligned}
}{Objetivo intermedio: agregados monetarios. Diseño tradicional de la Política Monetaria}

Un ejemplo claro sería: (i) alterar el coeficiente de caja; (ii) la oferta monetaria experimentará una variación más que proporcional; y, (iii) se conseguirá en última instancia disminuir o aumentar la inflación en base a nuestro objetivo. Ahora bien, esto requiere suponer que existe una relación estructural entre los objetivos últimos y los instrumentos utilizaods, lo cual no es siempre cierto:
\begin{lnum}
    \item \textbf{Problemas de identificación}. Analizar la evolución de los objetivos finales (variables) no es adecuado para analizar la evolución e impacto de la Política Monetaria ya que, como ha sugerido la evidencia empírica, la cuantifican del impacto es difícil, experimenta retardos temporales, etc;
    \item \textbf{A nivel operativo}. Además, a nivel operativo esto supone un problema añadido: ¿cuándo, cómo y cuánto reorientas la Política Monetaria? La incertidumbre e información imperfecta son cualidades inherentes del sistema, por tanto, la la valoración de las actuaciones resulta extremadamente compleja si se fundamenta sobre indicadores que emanan del sistema.
\end{lnum}

\subsubsection{Experiencia histórica: \textit{Money Targeting Strategy}}
Esta metodología de diseño fue implementada en muchas economías durante los años 70's. En particular, destacan los ejemplos de Reino Unido, canadá, Alemania y Suiza. 

A modo de síntesis:
\begin{lnum}
    \item \textbf{Objetivo inflación}. El objetivo último de sus Autoridades era el control de la inflación;
    \item \textbf{Instrumentos u objetivos intermedios: agregados monetarios}. La conducción de la Política Monetaria se instrumentalizaba como función de los agregados monetarios. Es decir, el objetivo intermedio de estas autoridades era el control del crecimiento de los agregados monetarios (lo que hemos evaluado), evitando en cualquier caso desviaciones grandes sobre los objetivos intermedios.
\end{lnum}

De hecho, por su instrumento este diseño pasó a ser conocido como \textit{Monetary Targeting Strategy}. Los resultados fueron heterogeneos:
\begin{la}
    \item \textbf{Fracaso.} Se podría decir, saalvando las distancias, que los gobiernos de Canadá, Estados Unidos y Reino Unido fracasaron en su objetivo de estabilización de la inflación. A pesar de mantenerse fieles a su objetivo intermedio (en particular, la no variación de la tasa de crecimiento de los agregados) no consiguieron estabilizar la inflación. Podría decirse que esto se vió fuertemente influenciado por la persecución de objetivos complementarios como la estabilidad del tipo de cambio y la, posible, inestabilidad estructural experimentada entre la tasa de crecimiento de los agregados monetarios y el PIB nominal (probablemente también por factores coyunturales, como por ejemplo un shock de oferta: crisis del petroleo);
    \item \textbf{Éxito}. Otros gobiernos, como Suiza y Alemania si alcanzaron su objetivo de estabilizar su inflación. Lo que resulta más revelador del éxito de estas autoridades es el énfasis que pusieron en la transparencia y comunicación de sus objetivos. Por ejemplo, el BundesBank se marcó unos objetivos de inflación a largo plazo a tipo de \textit{meta} y lo consiguió mediante el ajuste gradual a partir de metas a más corto plazo (medio plazo) a pesar de ser, o precisamente por ser, mucho más flexibles en lo que respecta a sus objetivos intermedios (se experimentaron desviaciones importantes en el crecimiento de los agregados monetarios).
\end{la}

\paragraph{Leccines aprendidas: Comunicación y transparencia}
¿Qué podemos aprender de esta experiencia histórica? Precisamente, el éxito de la estrategia seguida por dichas AM's se basó en la instrumentalización de los agregados monetarios como mecanismo de \textbf{comunicación} del objetivo de inflación que se habían fijado (no tanto como un objetivo intermedio confiando en relaciones estructurales perennes). 

Además, respecto a la inestabilidad percibida entre los agregados monetarios y el PIB nominal, muchos lo atribuyen a las sucesivas innovaciones financieras que se experimentaron a lo largo de esa y la siguiente década (inicios de los medios de pago electrónico) sumado a la inestabilidad financiera que sufrían los mercados por la incertidumbre generado por la crisis del petroleo (y contexto geopolítico global). Esta inestabilidad mostró:
\begin{lnum}
    \item Una alta volatilidad/varianza de la demanda de dinero, que hasta entonces se había considerado relativamente estable; y,
    \item Las dificultades de estimación/seguimiento de los agregados monetarios. Esto condujo, de hecho, a un arduo debate acerca la deseabilidad de instrumentar la Política Monetaria, y por tanto fijarse unos objetivos intermedios, basados en los agregados monetarios o en los tipos de interés. Dado que el objetivo de la Política Monetaria es reducir la varianza experimentada de la renta nominal, y que la varianza en la demanda de dinero es mayor que la de la renta (curva LM es más inestable y Curva IS es más estable), seria mejor articular la Política Monetaria como objetivo intermedio.
\end{lnum}

\begin{specialblock}{Instrumento deseable: POOLE (1970)}
    Esto último es precisamente el objeto de uno de los trabajos más célebres en relación con la Política Monetaria.

    En 1970, POOLE se realiza la siguiente pregunta ¿debería el banco central mantener los tipos de interés constantes o debería de mantener la oferta monetaria constante, permitiendo la libre fluctuación de los tipos de interés? Asumirá que el objetivo final de la política monetaria será estabilizar el PIB, por lo que la función objetivo a minimizar es la varianza del $\sigma_{PIB}=E[x_t]^2$.

    Bajo el contexto del modelo IS-LM, donde: (i) la curva IS representa el lado real de la economía, de tal fonna que existe una relación negativa entre PIB y el tipo de interés; (ii) la curva LM representa el lado monetario de la economía, donde la demanda de dinero depende positivamente del PIB y negativamente del tipo de interés; y (iii) ambas curvas estarán afectadas por shocks estocásticos ($u_t$ y $v_t$, respectivamente). Realiza las siguientes regresiones:
    \begin{la}
        \item Curva IS: $x_t=-\alpha i_t+u_t$;
        \item Curva LM: $m_t=x_t-c\;i_t+v_t$
    \end{la}

    El análisis de Poole pennite la representación gráfica de la volatilidad del PIB en función de si se fija $m$ o $i$ como instrumento operativo y, a su vez, si la economía vive shocks de naturaleza real o monetaria:

    \imagenfit{POOLE_1970.png}{Representación gráfica del análisis de POOLE}{Tema 3.A.35. Apuntes ICEX-CECO}

    Por tanto:
    \begin{la}
        \item \textbf{Shocks reales}. Bajo shocks de naturaleza real, representados por fluctuaciones de la curva IS, el control de la oferta monetaria permite estabilizar parcialmente la fluctuación del PIB ante el incremento de la demanda agregada: dada una oferta monetaria constante, el incremento de la demanda de dinero derivado del aumento de la actividad real se ajusta con un aumento del tipo de interés, lo que genera un crowding out parcial de la demanda de inversión. Por el contrario, mantener un tipo de interés constante implica que, ante las presiones al alza sobre los tipos derivadas del aumento de la demanda de dinero, el banco central las acomoda con una expansión monetaria, reflejada en una expansión de la curva LM. En conclusión, la elección del tipo de interés como instrumento de política monetaria amplifica los efectos de los shocks reales, mientras que el control de la oferta monetaria reduce su impacto; y,
        \item \textbf{Shocks monetarios}. Ante un shock monetario, como una caída en la demanda de dinero, el control de la oferta monetaria no consigue estabilizar la actividad económica. La caída en la demanda de dinero es compensada con una caída de los tipos de interés. lo que conduce a un incremento de la inversión, aumentando el PIB de la economía. Por el contrario, si la política monetaria mantiene constante el tipo de interés, la caída en la demanda de dinero es compensada con una contracción en la oferta monetaria, estabilizando por completo a la economía.
    \end{la}

    Dada una varianza constante de los shocks monetarios y reales ($\sigma_u^2=\sigma_v^2$), la instrumentación de la política monetaria a través del tipo de interés será más deseable cuanto más empinada sea la curva LM (derivada de una mayor sensibilidad de la demanda de dinero a variaciones en el PIB y/o a una menor sensibilidad a variaciones en el tipo de interés) y cuanto más plana sea la curva IS (derivada de una mayor sensibilidad de la demanda de inversión ante variaciones en el tipo de interés).

    El correcto funcionamiento del monetary targeting depende de un elemento clave: que la velocidad del dinero sea estable. Si ésta es inestable (lo que es equivalente a la materialización de shocks monetarios), la relación existente entre los agregados monetarios y el nivel de precios se disolverá. Fruto de la innovación financiera durante los años 70, la velocidad del dinero comenzó a fluctuar de manera inestable. Una vez fue claro a principios de los años 80 que la relación entre oferta de dinero y renta ya no era fuerte, EEUU, Canadá y Reino Unido decidieron abandonar de manera formal el monetary targeting. En palabras del Gobernador del Banco de Canadá en aquel momento, Gerald Bouey: \textit{Nosotros no abandonamos la estrategia del monetary targeting, sino que ella nos abandonó a nosotros}.
\end{specialblock}

\subsection{Diseño contemporaneo de la Política Monetaria: \textit{Inflation Targeting Estrategy}}
Por tanto, como vemos, las virtudes y deficiencias de la estretagia anterior dejaron en evidencia una cuestión esencial a la que poca importancia se le había dado, más allá de la fuerte influencia en el ámbito académico: las \textbf{expectativas} y la \textbf{confianza} son esenciales en la transmisión de la Política Monetaria.

A partir de los años 90's, el consenso era abrumador: las actuaciones de Política Monetaria debían diseñarse como teniendo en cuenta a los agentes de la economía. Las decisiones por tanto se vieron reducidas a una única variable de reacción: $PM=f(\text{Desviaciones de }\pi^e\;\;\text{respecto a }\pi^{Obj})$. Las características que emergieron de esta nueva forma de conducir/diseñar la Política Monetaria son:
\begin{lnum}
    \item \textbf{Compromiso}. La importancia del compromiso en alcanzar objetivos finales, en particular la estabilidad de precios (objetivo de control de la inflación); y,
    \item \textbf{Transparencia}. Los objetivos debían comunicarse al público, tanto los objetivos a medio plazo como los objetivos a largo plazo.
\end{lnum}

En resumidas cuentas, la actuación de la AM debía centrarse en garantizar la confianza de los individuos en el sistema y sobre ésta misma, de forma que los objetivos fijados fuesen creíbles y condicionaran la formación de las expectativas de los agentes. Además, este tipo de diseño fue conocido como \textit{inflation targeting}, tanto por su objetivo como por su claridad. 

\subsubsection{Experiencia histórica}
A partir de la experiencia de los países que han adoptado esta estrategia, podemos ver como los resultados han sido existosos por lo que respecta a:
\begin{lnum}
    \item \textbf{Anclaje de las expectativas}. La elección de un objetivo único de inflación (\textit{inflation targeting}), bien sea fijo o en bandas de fluctuación ha conseguido aumentar la credibilidad de las Autoridades Monetarias alnacanzando los objetivos fijados;
    \item \textbf{Acomodaticia}. El establecimiento de objetivos a medio y largo plazo ha tenido dos implicaciones muy importantes: (i) la Política Monetaria ha conseguido acomodar regularmente ciertos shocks (debido a su flexibilidad); y, (ii) la posibilidad de atender a otros objetivos secundarios pero importantes a través de los objetivos fijados a medio plazo (\textit{Flexible Inflation Targeting});
    \item \textbf{Disciplinador o no-discrecionalidad}. La transparencia y comunicación exigida ha tenido un efecto disciplinador en la actuación de las Autoridades Monetarias. En particular, una reducción de la incertidumbre generada por las actuaciones discreccionales y un mayor grado de independencia del Banco Central; y,
    \item \textbf{Indicadores globales homogeneos}. No debemos olvidar que la homogeneización del diseño de Política Monetaria y la importancia que en las últimas décadas se le ha otorgado ha permitido ampliar el abanico de indicadores de los que disponen las AM para realizar su labor. Esto, indudablemente, ha facilitado y facilita la toma de decisiones al anticipar presiones inflacionistas pero también ha aportado grandes posibilidades a la ciencia económica en su conjunto al agrandar la disponibilidad de datos.
\end{lnum}

\paragraph{Limitaciones: Mejoras}
Sin embargo, cabe destacar también las limitaciones que esta mantiene, como por ejemplo los retardos en la actuación de la Política Monetaria a la hora de incidir en las variables últimas objetivo (aunque pueden considerarse inherentes al funcionamiento del sistema por la existencia de fricciones, sí cabe destacar que las innovaciones financieras deben reducir éstas). 

\subsection{Los objetivos reales: Debates y controversias persistentes}
Ahora bien, de todo el análisis histórico expuesto que nos ha permitido entender cómo se diseña la Política Monetaria y por qué se diseña de esta forma, sacamos también una conclusión controvertida: ¿no se espera de la Autoridad Monetaria una labor de estabilización real o colchón, por ejemplo, contracíclico?

Dijimos al inicio de nuestra exposición que obviavamos las consideraciones relativas a las Políticas Monetarias \textit{activistas}. ¿Pero es esto razonable? ¿Debe estar la Política Monetaria al margen de, por ejemplo, la estabilidad del empleo? De hecho depende, las Autoridades Monetarias del Eurosistema siguen un estricto compromiso de control de la inflación pero, por ejemplo, la FED tiene un compromiso doble: el control de la inflación y la estabilidad del empleo. 

Entramos aquí en el corazón de una de las contraversias más duraderas de la ciencia económica: el papel del dinero (y en consecuencia, de la Autoridad Monetaria) en la economía real. Sin ánimo de entrar en los argumentos de una u otra posición, sí cabe mencionar algunos resultados estrechamente relacionados con el diseño de la Política Monetaria.

\subsubsection{¿La dualidad imposible?: Crecimiento y estabilidad}
En primer lugar, podemos preguntarnos qué es mas deseable, ¿perseguir múltiples o un único objetivo? Mucho se ha teorizado al respecto, siendo una de las aportaciones más relevantes al respecto el modelo de KYDLAND y PRESCOTT (1977).

Estos autores demuestran que cuando se persiguen diferentes objetivos pueden surgir problemas de inconsistencia dinámica de la Política Monetaria, es decir, que una Política óptima \textit{ex ante} sea subóptima evaluado en el momento en el que se va a realizar (\textit{ex post}). Estos autores plantean un sencillo modelo en el que la Función Objetivo de la Autoridad Monetaria contempla tanto el nivel de desempleo como la estabilidad de precios (inflación): (i) a través de la Función de Pérdida Social, la Autoridad Monetaria desea minimizar tanto el nivel de desempleo como de inflación como objetivo, ya que ambos suponen un coste en términos sociales; y, por otro lado, (ii) su actuación se ve limitada, ya que existe un trade-off establecido por la relación a corto plazo representada por la Curva de PHILLIPS, atendiendo a que los agentes forman sus expectativas de precios en base a su conocimiento previo (HER).

\eqblock{
\begin{aligned}
    \boxed{\begin{aligned}
        &\min_{u_t,\pi_t}\overbrace{L=L(\pi,u)}^{\text{\scriptsize Función de Pérdida Social}}\\[1em]
        &\text{s.a.}\quad \underbrace{u=u_N+\frac{1}{\alpha}(\pi-\pi^e)}_{\text{\scriptsize Curva de PHILLIPS}}
    \end{aligned}}
\end{aligned}
}{Modelo KYDLAND y PRESCOTT: Objetivos Múltiples. Objetivos y controversias de Política Monetaria}

Atendiendo al grado de importancia que le otorgue la autoridad monetaria a uno u otro objetivo, tendremos unas curvas de indiferencia con más o menos pendiente (siempre negativa), pero supongamos que le otorga el mismo peso a ambos objetivos:

\imagenfit{ID_PuntoInicial_KYDL_PRESCOTT.png}{Situación inicial en términos de objetivos de la Autoridad Monetaria}{Elaboración propia}

Como vemos, la economía se encuentra en un punto de equillibrio ($A$) donde la tasa de variación de los precios es nula (ha alcanzado la estabilidad de precios) y el nivel de desempleo se encuentra en su nivel natural (determinado por condiciones reales de la economía, como tecnología, producción, etc.). Ahora bien, la Autoridad Monetaria puede alterar el aumentar el nivel de empleo aceptando un mayor grado de inflación sin perder Utilidad Total (es más óptimo), es decir, manteniendose en la misma curva de indiferencia y, por tanto, en el mismo grado de (des)utilidad social. Esto se realizaría conduciendo la economía, la inflación, al nivel de intersección de esta curva con la restricción de maniobra (restricción de optimización):

\imagenfit{ID_SubOptimo_KYD_PRESCOTT.png}{(Izq.) Posibilidad de explotar la Curva de Phillips a corto plazo por parte de la AM | (Dcha.) Anticipación de los agentes y nuevo equilibrio subóptimo}{Elaboración propia}

Por tanto, como vemos, la AM tiene margen de maniobra en este punto inicial de forma que, explotando la Curva de PHILLIPS a corto plazo se alcanza un mayor nivel de empleo, soportando más inflación pero manteniendo la misma (des)utilidad social. Sin embargo, como los agentes anticipan este movimiento por parte de la AM (inflación esperada), el nuevo equilibrio se alcanzará en el punto ($A'$). Dando lugar a un mismo nivel de empleo pero con mayor nivel de inflación, aumentando la (des)utilidad social y por tanto, empeorando la situación. Se alcanzaría un Equilibrio Sub-óptimo por el problema de inconsistencia dinámica que se presenta. 

De hecho, aunque pueda conseguirse lo que se conoce como \textbf{sorpresa monetaria}, el punto $B$ alcanzado no sería estable, puesto que los agentes revisarían sus expectativas de inflación y el equilibrio retornaría al punto estable, siendo éste el único Equilibrio de NASH. Además, cabe destacar que, en este punto ($A'$), la Autoridad Monetaria se \textbf{ata las manos} puesto que no es capaz de explotar ninguna relación que le otorgue la Curva de PHILLIPS a corto plazo, se trata de la intersección de la curva de desempleo natural (estructural), la Curva de PHILLIPS a corto plazo y la Función de Pérdida Social.

Las dos conclusiones más importantes de este modelo son:
\begin{lnum}
    \item \textbf{La credibilidad}. En primer lugar, si la Autoridad Monetaria desea mantener una mínima posibilidad de tener algún tipo de injerecia sobre las variables reales de la economía, debe ser realmente creíble en el objetivo de control de precios; puesto que, como se ha visto, la única forma de trasladar una Política Monetaria expansiva aumentando el nivel de empleo es que los agentes no anticipen este incremento, que experimenten sorpresa monetaria. Aunque este no sea estable; y,
    \item \textbf{Objetivo único}. Independientemente de la posibilidad de explotar la Curva de PHILLIPS a corto plazo, los autores resaltan la importancia de perseguir un único objetivo cuantificable, en particular, la estabilidad de precios (\textit{inflation targeting}). Cualquier otra intervención u objetivo, conduciría a la economía a los mismos resultados en términos reales pero con un mayor nivel de inflación (menoscabando también la credibilidad de la AM)
\end{lnum}

\subsubsection{La Política Monetaria discrecional: ventajas e inconvenientes}
Habida cuenta de lo expuesto, podríamos preguntarnos cómo y cuál es la mejor manera de perseguir este objetivo (lo mismo aplicaría para múltiples objetivos). En particular, ¿debemos optar por una Política Monetaria discrecional o reglada en persecución de los objetivos?

La ruptura del compromiso de estabilidad de precios para alcanzar un mayor nivel de empleo, enunciado en el modelo anterior, no solo pone de relieve la bondad de perseguir un único objetivo (\textit{inflation targeting}, por ejemplo), sino también la importancia de la existencia de reglas que incrementen la credibilidad de la AM, por dos razones: (i) una Política Monetaria reglada ata las manos a la Autoridad Monetaria, lo que conduce a un aumento de credibilidad por parte de los agentes y, por tanto, es más sencillo alcanzar el objetivo perseguido (aplicado al caso anterior, lo que conocemos como Equilibrio de RAMSEY: $u_t=u_n\;\;,\;\;\pi=\pi^e=0$); y, por otro lado, (ii) un mayor grado de discrecionalidad en la Política Monetaria aumenta los incentivos de la AM para desviarse, de manera que se alcance, en este caso, el Equilibrio de NASH.

\textbf{As Taylor (1993) noted, “If there is anything about which modern macroeconomics is clear however—and on which there is substantial consensus—it is that policy rules have major advantages over discretion in improving economic performance”}

Desde otro punto de vista, (GALI, CLARIDA y GERTLER, 1999) algunos autores defienden que los resultados obtenidos son similares siguiendo una Política Monetaria discrecional o reglada, sobre en términos de minimizar las pérdidas sociales. Ahora bien, los mismos autores establecen que éstos serán similares siempre y cuando:
\begin{lnum}
    \item \textbf{PM reglada}. Si se opta por una Política Monetaria reglada, la fijación del objetivo de precios tenga especialmente en cuenta las expectativas que los agentes forman respecto a las condiciones económicas futuras. Lo que a su vez permitiría a la Autoridad Monetaria enfrentarse a un \textit{trade-off} ($\pi,u_t$) mejorado (reforzado); o, por otro lado,
    \item \textbf{PM discrecional}. Si se opta por una Política Monetaria discrecional, la Autoridad Monetaria asigne (o pondere) al control de la inflación un peso mayor en su actucación, lo que permitiría reducir el sesgo inflacionista aunque se mantenga una Política Monetaria discrecional.
\end{lnum}

\subsubsection{Tipología de Reglas}
Por último, antes de presentar los instrumentos de Política Monetaria, cabe hacer un breve repaso de la tipología de reglas que podemos considerar en Política Monetaria. En particular, el debate se centra en tres cualidades esencialmente.

\paragraph{Reglas Sencillas vs. Optimización}
¿Deben las reglas presentarse de una manera sencilla o deben las reglas seguir procesos de optimización continuos? ¿Qué diferencia hay entre estas opciones?
\begin{lnum}
    \item La primera alternativa sería vincular, por ejemplo, la variable intermedia (sobre la que incide la Política Monetaria) a un conjunto reducido de indicadores económicos; por otro lado,
    \item La segunda alternativa se basaría en considerar un mayor abanico de variables en un marco modelizado, por ejemplo, un modelo económetrico de forma que podamos estudiar las reacciones entre estas y de este forma, escoger nuestra \textbf{regla óptima}.
\end{lnum}

Aunque la diferencia es sustancial, la evidencia empírica sugiere que es más recomendable seguir reglas sencillas. Estas son más robustas y pueden utilizarse eficazmente en contextos más heterogeneos (\textit{frameworks} diversos).

\paragraph{Flexibles vs. Rígidas}
Si es mejor establecer un conjunto de reglas sencillas, ¿deben ser estas flexibles o rígidas?

Un conjunto de reglas flexibles tendrían en cuenta la posibilidad de comodarse mejor a coyunturas económicas, reaccionando ante shocks inesperados. Por ejemplo, aquí encontraríamos el diseño de \textit{inflation-targeting} mediante bandas de fluctuación y objetivos múltiples pero priorizando la estabilidad de precios: la FED.

Por el contrario, un conjunto de reglas rígidas sería, entre muchas otras, la Regla de FRIEDMANN respecto a la tasa continua y constante de incremento de los agregados monetarios, o un \textit{inflation targeting} más definido: el BCE.

\subsubsection{¿Reglas sobre qué? Objetivos vs. Instrumentos}
Por último, podríamos preguntarnos ¿qué es más óptimo reglar? Es decir, ¿qué debería mantener como objetivo la Autoridad Monetaria? En particular, ¿deberíamos fijar reglas sobre los objetivos a alcanzar o sobre los instrumentos a utilizar?

Si nos acordamos de las diferencias en el diseño de la Política Monetaria actual respecto a la seguida en el pasado, podría decirse que ese debate quedo zanjado al evidenciarse las limitaciones que presentaba la fijación de objetivos intermedios para controlar los objetivos finales. Sin embargo, esto es a la vez cierto y erroneo. La fijación de objetivos últimos condujo a lo que se cononce como la \textbf{Gran Moderación} (ca. 1980-2007), donde la actuación de la Autoridad Monetaria en base a la fijación de objetivos finales mantuvo la inflación bajo control y un crecimiento relativamente estable en la mayoría de economías desarrolladas. 

\imagenfit{NGDP_Trends.png}{Two NGDP Periods}{\textit{Facts, Fears and Functionality of NGDP Targeting: A Guide to a Popular framework for Monetary Policy}. BEECKWORTH (2019)}

Ahora bien, la mayor aportación de este enfoque no fue poner de relieve las virtudes de la fijación de objetivos, sino trasladar el papel de las expectativas a una posición preeminente en la actuación de la Política Monetaria. 

Durante las últimas dos décadas, las AM's han enfrentado serios problemas a la hora de conducir la economía a sus objetivos mediante los instrumetos tradicionales, que aunque serán tratados en el apartado siguiente, parecen no ser suficientes para estimular la economía. Por ello es relevante preguntarse sobre qué debería reglar/proponerse la AM y cómo esto puede ayudarnos a salir de problemas como, por ejemplo, el \textit{Effective-Lower-Bound}.

Sin entrar en disquisiciones teóricas sobre los objetivos finales y su bondad, por haberse ya expuesto los virtudes del \textit{inflation targeting} y como esto ha llevado a la \textbf{Gran Moderación}; pasemos directamente a analizar sus alternativas más relevantes y cómo estas podrían ayudarnos en la actualidad. Distinguiremos dos propuestas que han contado con importante respaldo en la literatura, y con aplicación parcial por parte de algunas economías desarrolladas:

\paragraph{Tipo de Interés: \textit{Price-Level Targeting}}
El Price-Level-Targeting (PLT) es una de las principales alternativas al \textit{inflation targeting} (IT), según la cual se podrían evitar algunas de las mayores limitaciones asociadas a la aplicación del IT en un contexto de Effective Lower Bound (ELB). 

La principal diferencia entre el PLT y el IT es que, mientras que el IT busca alcanzar un nivel de inflación promedio, el PLT busca seguir cierta trayectoria del nivel de precios establecida de antemano (objetivo o regla). Por tanto, ante shocks inflacionarios en períodos anteriores, el IT los tomará como errores del pasado y seguirá intentando alcanzar el objetivo de inflación, mientras que el PLT buscará corregirlos para volver a converger con la senda de precios objetivo. La diferencia entre ambos puede observarse en el siguiente gráfico:

\imagenfit{PLT_vs_IT.png}{Inflation Targeting and Price-Level-Targeting compared}{\href{https://cepr.org/voxeu/columns/inflation-targeting-vs-price-level-targeting-new-survey-theory-and-empirics}{\textit{Inflation Targeting vs Price Level Targeting: A new survey of theory and empirics}. MINFORD y HATCHER (2014)}}

Bajo un régimen de PLT, el banco central está obligado a corregir los shocks inflacionarios pasados.\footnote{En consecuencia, este régimen es \textit{history dependent} (WOODFORD, 2003).
} Por ejemplo, consideremos que la economía sufre un shock negativo sobre la demanda agregada, lo que obliga al banco central a reducir los tipos de interés hasta el ELB (digamos, 0\%). Bajo un régimen de IT, dado que las expectativas de inflación tras el shock siguen siendo inicialmente las del target (supongamos, 2\%), familias y empresas reconocerán que la política monetaria ya ha alcanzado sus límites y que ya no es capaz de estimular más la economía, lo que podría llevar incluso a reducir las expectativas de inflación, incrementando el tipo de interés real y deprimiendo aún más a la economía: se avecina una larga recesión. Por el contrario, en un régimen de PLT creíble, las expectativas de inflación no permanecen ancladas al 2\%. 

¿Por qué? Aquí está la esencia del PLT: ante un shock de demanda negativo que genera una inflación presente inferior al target, los agentes esperarán una política monetaria futura aún más expansiva que genere una inflación presente superior al target del 2\%. Las mayores expectativas de inflación, dado un tipo de interés nominal fijo, reducirán los tipos de interés reales, estimulando la demanda agregada y la producción, permitiendo así escapar del ELB.\footnote{%
    Esta argumentación lógica ha sido posteriormente confirmada en modelos de la NEK, por ejemplo, EGGTERSSON y WOODFORD (2003) o NAKOV (2008). En estos modelos, las pérdidas de bienestar en contextos de ELB son significativamente superiores bajo IT que bajo PLT.
}

\paragraph{Agregados Monetarios: \textit{Nominal GDP targeting}}
Otra alternativa al IT es el \textit{Nominal GDP targeting} (NGDPT), según el cual los bancos centrales tienen como objetivo un crecimiento estable del gasto nominal que se produce en la economía, que es igual a la renta nominal. 

Como sabemos, dicha renta nominal ($P\cdot Y$) se descompone en dos elementos: el nivel de precios y la renta real. Por tanto, el NGDPT ancla el tamaño nominal de la economía y su ritmo de crecimiento, lo que a su vez ofrece un ancla de referencia con respecto a la tasa de crecimiento de precios y salarios. Al igual que el PLT, el NGDP ofrece un mecanismo para escapar de la ELB ante shocks que produzcan caídas en el nivel de actividad y de inflación. Al encontrarse las expectativas ancladas con respecto a una senda de evolución del PIB nominal, los agentes esperarán una política monetaria más expansiva en el período futuro que permita compensar el shock actual:

\imagenfit{NGDPT.png}{A Nominal GDP Level Target}{\textit{Facts, Fears and Functionality of NGDP Targeting: A Guide to a Popular framework for Monetary Policy}. BEECKWORTH (2019)}

Por otra parte, el NGDPT puede ser entendido como un sistema target de la oferta monetaria de dinero, corregido por los cambios en su velocidad. Dado que $P\cdot Y=M\cdot V$, al establecer como objetivo cierto gasto nominal, este régimen también permite neutralizar cambios en la velocidad del dinero, que se verán compensados por cambios inversamente proporcionales en la oferta monetaria. Otra de las virtudes del NGDPT es que es ofrece un marco de respuesta robusto tanto a shocks de oferta como a shocks de demanda. 

\imagenfit{NGDPT_SupplyShocks.png}{Nominal GDP Targeting and Supply Shocks}{\textit{Facts, Fears and Functionality of NGDP Targeting: A Guide to a Popular framework for Monetary Policy}. BEECKWORTH (2019)}

Bajo el IT, ante un shock de oferta, para luchar contra la inflación, el banco central puede generar una subida excesiva de los tipos de interés que acabe provocando una intensa recesión económica. No obstante, en este tipo de situaciones, el banco central debería de acomodar el incremento en los precios debido a factores de oferta (como incrementos en los precios relativos o cambios en la capacidad productiva de la economía) y centrarse exclusivamente en las presiones inflacionistas motivadas por los factores de demanda. El NGDPT hace automáticamente esto al centrarse en un crecimiento constante del PIB nominal. 

Los defensores del NGDPT también ensalzan es su capacidad para limitar la inestabilidad en los mercados financieros. Ello se debe a que la fonnación de burbujas está fuertemente correlacionada con crecimientos del PIB nominal por encima de la media. Por ejemplo, en el período 2002-2004, en EEUU, dado el shock positivo de productividad que se estaba viviendo, la FED decidió mantener los tipos de interés por debajo del nivel óptimo, debido a sus preocupaciones por una posible desinflación, aún a pesar de que durante este período el PIB nominal estaba creciendo por encima de la media. El NGDPT evita esto al tener como objetivo un crecimiento constante del PIB nominal.


\clearpage
\section{Instrumentos de la Política Monetaria}
\puntosclave{
    El banco central no controla directamente la oferta monetaria, pero sí puede afectar a sus determinantes a través de sus instrumentos para así poder influir finalmente sobre el nivel de precios.
    
    La autoridad monetaria ha de decidir su instrumento de conducción: la base monetaria o el tipo de interés. El análisis de POOLE nos permite identificar qué instrumento será el óptimo en función de la magnitud de los shocks reales y monetarios.
    
    En la práctica, los Bancos Centrales intervendrán en el mercado monetario a través de diferentes instrumentos que se pueden clasificar en torno a tres categorías: intervención vía cantidades (incidencia sobre los agregados monetarios), vía precios (incidencia sobre los tipos de interés) o, vía comunicación (incidencia sobre las expectativas.
    
    La adopción de las CBDC supondría crear un nuevo componente en el agregado monetario $M_0$ (base monetaria): se sumaría así al efectivo y las reservas, lo que daría lugar a la aparición de un nuevo instrumento a través del cual implementar la política monetaria
}
Ahora que hemos visto los diferentes enfoques en el diseño que puede seguir la Política Monetaria, pasemos a analizar los instrumentos que la Autoridad Monetaria posee para conseguir sus objetivos.

Independientemente de la estrategia de Política Monetaria perseguida, ésta se instrumenta, sobretodo, mediante modificaciones de liquidez del Banco Central a los Bancos Comerciales (o intermediarios de crédito/financiero).

\subsection{Instrumentos Primarios (cuantitativos)}
Los instrumentos primarios, o cuantitativos, son principalmente tres.

\subsubsection{Políticas de Reservas Mínimas (coeficiente de caja)}
El primer instrumento es el coeficiente de reservas mínimas, que obliga a mantener un porcentaje determinado sobre las reservas mantenidas en los activos de los Bancos Comerciales (que pueden o no estar depositados en el Banco Central).

Por tanto, a través del coeficiente de reservas mínimas y su proyección sobre el primer agregado en circulación ($M_1$), hemos visto como el Banco Central incide en la Oferta Monetaria. Aumentando (disminuyendo) el coeficiente de reservas mínimas, el banco central disminuye (aumenta) la Oferta Monetaria en la Economía. 

El objetivo perseguido, principalmente, a través de este instrumento es el de la estabilización de los tipos de interés nominal, ya que, a modo de ejemplo, un incremento del coeficiente legal de reservas, disminuiríra la Oferta Monetaria y esto incidiría negativamente sobre las posibilidades de concesión de crédito y la creación de dinero bancario.

\subsubsection{Política de Crédito y Redescuento}
Otro de los instrumentos utilizados por las Autoridades Monetarias, principalmente en el corto plazo, son los instrumentos vía precios. 

Se trata en esencia de la concesión de créditos a corto plazo por parte del banco central a las entidades de crédito y del depósito de excesos de liquidez de las entidades de crédito en el banco central. Se distinguen dos tipos de facilidades:
\begin{la}
    \item \textbf{Facilidades de crédito}. La facilidad de crédito permite a las entidades obtener liquidez del Banco Central a un día a tipos sensiblemente superiores a los de mercado, contra activos de garantía. Permite, por tanto, establecer un techo a los movimientos de los tipos de interés en los mercados monetarios; y,
    \item \textbf{Facilidades de depósito}. La facilidad de depósitos permite a las entidades de crédito depositar sus excedentes de tesorería en el banco central. Estos depósitos están remunerados a tipos inferiores a los de mercado. Permite, por tanto, establecer un suelo a los movimientos de los tipos de interés en los mercados monetarios. 
\end{la}

El objetivo es el de proporcionar o absorber liquidez a un día, señalar la orientación de la política monetaria y controlar los tipos de interés del mercado monetario.\footnote{
\textbf{Tipo de interés del mercado monetario y Tipo de interés de mercado}. Al igual que el banco central no dispone de un control perfecto sobre la oferta de dinero, tampoco controla con precisión el precio de equilibrio del mercado monetario: el tipo de interés de mercado. Partimos de la constatación de que el sistema financiero está conformado por diferentes agentes: bancos, inversores, prestatarios y ahorradores. Cada uno tiene sus propias preferencias, aversión al riesgo y horizontes de inversión. Estos factores influyen colectivamente en la oferta y la demanda de fondos prestables, que, a su vez, determinan los tipos de interés del mercado. Por otra parte, cabe señalar que no existe un único tipo de interés de mercado, sino un continuo de ellos, que varían por un lado según los distintos tipos de préstamos e inversiones, reflejando la calidad crediticia y los perfiles de riesgo subyacentes, y por otro en función del plazo del instrumento financiero. Los Bancos Centrales fijan los tipos de interés oficiales, que influyen principalmente sobre los tipos de interés a muy corto plazo, pero no controlan las primas de riesgo asociadas a los distintos tipos de préstamos o inversiones ni las primas temporales (que vienen detenninadas por la relación entre preferencia por el consumo presente frente el futuro -$\rho$- así como por las expectativas). Son los agentes del mercado los que evalúan la solvencia de los prestatarios, las expectativas futuras de las condiciones del mercado monetario y la incertidumbre asociada a las mismas y, en función de estos factores, ajustan los tipos de interés para compensar el riesgo asumido.
}

\paragraph{Pasillo Interbancario}
Varios bancos centrales en el mundo (Australia, Canadá, Nueva Zelanda, Suecia o Suiza) utilizan un sistema de pasillo interbancario para controlar los tipos de interés. 
    
En este sistema, a través de las facilidades de crédito y de depósito, el banco central crea un pasillo por el que siempre se moverá el tipo de interés interbancario. Este sistema le permite tener un mayor control al banco central de las fluctuaciones de los tipos, así como alterar el tipo de interés de mercado sin alterar el nivel de reservas (WOODFORD, 2001c). Para ilustrar simplemente qué es un pasillo interbancario podemos concebirlo como que: (i) el banco central paga a las reservas depositadas en él un tipo de interés igual a $i^*-s$ y cobra un tipo de interés al préstamo de reservas igual a $i^*+s$. En consecuencia, en este sistema ningún banco pedirá prestado en el mercado interbancario privado a un tipo de interés superior a $i^*+s$, de igual forma que nunca se prestará a un tipo de interés inferior a $i^*-s$. El tipo de interés interbancario siempre se encontrará en el siguiente intervalo: $i^*-s\leq i\leq i^*+s$.
    
\imagenfit{Pasillo_Interbancario.png}{Tipos de interés Oficiales (BCE) y el Tipo de interés Interbancario}{Banco de España (2023)}

El objetivo es, por tanto, acotar el comportamiento interbancario (no del mercado) en un entorno \textit{óptimo} ya que esto incide indirectamente en los tipos de interés del equilibrio de mercado (independientemente de la cantidad, número, etc. que haya). ¿Por qué? Porque en aras de buscar rentabilidad, se espera que el comportamiento de los bancos se acerque lo máximo posible al comportamiento que de éstos espera el Banco Central: (i) se prestarán entre ellos a un tipo ligeramente superior a la rentabilidad ofrecida por el Banco Central por sus depósitos; y, (ii) el tipo de interés de mercado se fijará en el umbral que quede entre el tipo de interés interbancario/facilidad de depósito y el tipo de interés de facilidad de crédito (si estos fueran inferiores, habría \textit{infinita garantía/rentabilidad} del sector (podrían solicitar préstamos, contra garantía, y \textit{cobrar} estos a un interés superior)).

\subsubsection{Operaciones de Mercado}
Ahora bien, estos instrumentos de control cuantitativo tradicionales han experimentado ciertas limitaciones a lo largo de los años, el primer ejemplo claro al respecto lo encontramos en la economía Japonesa a finales de la década de los 90's, fenómeno que se extendió a la gran mayoría de economías desarrollados a raíza de la crisis financiera de 2007. Pero, ¿qué pasó?

Dijimos en apartados anteriores como el \textit{inflation targeting} había logrado el objetivo de moderar la volatilidad de la inflación del periodo anterior. No obstante, el firme compromiso adquirido por los Bancos Centrales, podríamos hablar de una reputación de enemigo último de la inflación, y los instrumentos comentados han sido poco eficaces a la hora de combatir expectativas inflacionistas ancladas por debajo de los objetivos óptimos. Estas expectativas ancladas a la baja han sido acompañados se intentaron remediar mediante el establecimiento de coeficiente mínimos y tipos de interés cercanos a cero, sin que ninguna de estas medidas surtiera el efecto deseado de estimular la economía e incrementar la inflación. 

Por ello, en las últimas décadas se han buscado instrumentos alternativos de Política Monetaria que, sumado con los debates sobre las alternativas viables respecto a la fijación de objetivos, pudiesen ayudar a este respecto. En particular, han tomado especial relevancia las operaciones de mercado abierto realizadas por parte de los Bancos Centrales. De entre las posibles, destacamos en particular las siguientes.

\paragraph{Operaciones de Mercado Abierto}

\textbf{Nota.-} Mirar esta serie de videos para acompañar lo que se presenta si no queda claro: \href{https://www.khanacademy.org/economics-finance-domain/core-finance/money-and-banking/federal-reserve/v/fed-open-market-operations}{\textit{Fed Open Market Operations: QE and other operations}. KHAN Academy}

Las Operaciones de Mercado Abierto forman parte del abanico \textit{tradicional} de la Política Monetaria. Estas consisten en la compra y venta, por parte del Banco Central, de valores denominados en moneda nacional o extranjera, la cual determina también sus condiciones. 

Suelen efectuarse a precios de mercado y mediante subastas a las que sólo pueden acudir ciertas entidades autorizadas. Distinguimos dos tipos fundamentales de operaciones de mercado abierto: 
\begin{la}
    \item \textbf{Operaciones regulares}. Se ajustan a un calendario prefijado y constituyen la pieza estratégica de la intervención/instrumentación monetaria; y,
    \item \textbf{Operaciones de fine tuning}. Regula de forma más puntual las condiciones monetarias sobre la base de una mayor frecuencia y/o un plazo más corto.
\end{la}

Los objetivos principales son controlar los tipos de interés del mercado monetario, gestionar la liquidez del mercado y señalizar la orientación de la política monetaria.\footnote{
Esta política tiene efectos tanto sobre la cantidad de dinero, vía variación en la base monetaria, como sobre los tipos de interés, vía variación en el volumen de crédito.
}

\paragraph{\textit{Quantitative Easing} (QE)}
No obstante, el instrumento que ha suscitado más debate en los últimos años se conoce como \textit{Quantitative Easing} (QE). 

El término QE ha sido utilizado para describir diferentes programas. En este apartado nos vamos a ceñir a la siguiente definición: QE son programas de compra a gran escala de títulos a largo plazo, realizados por el Banco Central con el objetivo de reducir los tipos de interés a largo plazo, relajar las condiciones financieras y, en última instancia, alcanzar los objetivos macroeconómicos de pleno empleo o estabilidad de precios. 

Muchos economistas han criticado el QE tanto por ser el potencial causante de elevados niveles de inflación como de ser completamente inefectivo. No obstante, tal y como dijo BERNANKE en 2014, "\textit{El problema con QE es que funciona en la práctica, pero no en la teoría}". 

Los programas de expansión cuantitativa tienen efectos reales sobre la economía a través de los siguientes canales: 
\begin{lnum}
    \item \textbf{El canal balance-cartera}. Las compras del banco central de títulos a largo plazo reducen la oferta de los mismos que son mantenidos por el sector privado. Ello incrementa sus precios y, consecuentemente, reduce sus tipos de interés. Además, debido a que las entidades financieras pasan a tener un exceso de liquidez en sus balances, buscarán reducirla a través de la adquisición de otros títulos a largo plazo. De nuevo, el precio de estos activos aumentará y el tipo de interés implícito se reducirá;
    \item \textbf{El efecto señalización}. En consonancia con los conocimientos adquiridos respecto al papel de las expectativas, se considera que el anuncio de programas de QE puede dar la fuerte señal de que las autoridades están fuertemente comprometidas con mantener una política expansiva y unos tipos bajos durante varios períodos. Ello se debe a que: (i) los programas de QE duran meses y posiblemente años, y la autoridad monetaria no va a llevar políticas contradictorias subiendo tipos por un lado y manteniendo programas de compras de valores por otro, por lo que los programas de QE en cierta medida atan las manos a la autoridad monetaria; y, (ii) los programas de QE son costosos para el Banco Central, en consecuencia, al anunciar un programa de QE, la autoridad monetaria está señalizando su compromiso con mantener un apoyo sostenido a la economía.\footnote{%
        Estos programas han sido criticados políticamente y pueden provocar pérdidas de capital en el balance del banco central.
    }
\end{lnum}

En definitiva, QE no deja de ser un swap de activos financieros, donde activos menos líquidos (los títulos a largo plazo) son intercambiados por el activo más líquido en una economía, las reservas bancarias.

\paragraph{\textit{Yield Curve Control} (YCC)}
Este instrumento está conformado por el anuncio público y la ejecución de un programa de compra de activos financieros con un período de vencimiento concreto con el objetivo de garantizar que los rendimientos de dichos activos se mantengan en cierto nivel deseado. Dado que los precios de los bonos están inversamente relacionados con sus rendimientos, su compra conduce a un incremento de su precio, lo que empuja los tipos de interés hacia abajo. 

La principal diferencia con el QE es que, mientras que el énfasis de éste se encuentra en el anuncio de compra de cierta cantidad de activos financieros, el del YCC se encuentra en el compromiso de mantener el rendimiento de los bonos con cierto período de vencimiento al nivel anunciado de antemano, comprando y vendiendo tantos bonos como sea necesario. 

El Banco Central de Japón inició un programa de YCC en 20 1 6 para luchar contra las presiones deflacionistas a las que se venía enfrentando desde hacía años. En concreto, Japón, se comprometió a fijar el rendimiento de los bonos a 10 años al 0\%. Por otra parte, Australia y Nueva Zelanda también experimentaron con versiones de YCC como respuesta a la crisis del covid-19. Por ejemplo, Australia se fijó el rendimiento objetivo de los bonos a tres años al 0,25\%.

\subsection{Instrumentos de Política Financiera (cualitativos)}
Ahora que hemos visto los diferentes instrumentos cuantitativos de los que dispone la Autoridad Monetaria, pasemos a analizar otra serie de instrumentos de especial relevancia, tanto a raíz de la crisis financiera como por las limitaciones experimentadas respecto a los instrumentos tradicionales que hemos comentado: los instrumentos \textbf{cualitativos}.

\paragraph{Instrumentos Mercado Financieros}
En particular, destacan los relativos a la supervisión del funcionamiento del mercado financiero: los instrumentos de Política Financiera. Entre otros:
\begin{lnum}
    \item \textbf{Supervisión}. La supervisión de las entidades financieras, en lo relativo al riesgo ha resultado un elemento esencial a raíz de la crisis, con el fin de garantizar la estabilidad de estos mercados;
    \item \textbf{Control selectivo del crédito}. También destacan los intrumentos de control selectivo del crédito con el fin de garantizar que se dirija primordialmente a aquellos sector más relevantes;
    \item \textbf{Control de cambios}. La introducción de nuevas disposiciones legales y administrativas dirigidas a controlar o regular las transacciones (cobros y pagos) con el extranjero; y,
    \item \textbf{Persuasión moral}. Menos claros pero sí efectivos, los mecanismos de persuasión con el fin de convencer a las entidades de restringir préstamos a determinados grupos. 
\end{lnum}

\paragraph{\textit{Forward-Guidance}}
No obstante, el más relevante de los instrumentos cualitativos se conoce como \textit{forward-guidance} (FG). Este instrumento ha cobrado especial relevancia a lo largo de los últimos años y se puede definir como la comunicación realizada por \textit{policy makers} sobre sus expectativas de evolución tanto de la economía como de la política económica aplicada. 

La comunicación del Banco Central toma diferentes formas: discursos de las autoridades tras las reuniones de la directiva, la publicación de \textit{Monetary Policy Reports}, \textit{Inflation Reports}, etc. La intuición sobre por qué el FG funciona es que las condiciones financieras presentes dependen, además de los tipos de interés a corto plazo, de las expectativas del mercado sobre la evolución futura de los tipos. En consonancia (y a modo de acompañar) los nuevos instrumentos cuantitativos presentados, si el banco central consigue influenciar las expectativas del mercado de tal manera que se esperen menores tipos en el futuro, los tipos de interés a largo plazo caerán. 

De acuerdo con BERNANKE (2022), se distinguen dos tipos de FG:
\begin{la}
    \item \textbf{Guía délfica}. Su intención es únicamente la de informar y ayudar a los mercados a comprender las previsiones económicas de la autoridad monetaria, así como sus planes provisionales de política económica. En definitiva, es una previsión del banco central tanto de la coyuntura como de la política económica futura; y, por otro lado,
    \item \textbf{Guía odiséica}. Su intención es la de realizar una promesa o compromiso. Este tipo de guías son útiles principalmente cuando los tipos de interés se encuentran en el Effective Lower Bound. En este contexto, el banco central puede reducir aún más los tipos de interés a largo plazo persuadiendo a los mercados de que tienen la intención de mantener los tipos bajos durante un período más largo de lo que esperaban. Esto es lo que se ha conocido como políticas \textit{lower-for-longer}. En este sentido, debido a que la Guía odiseica es un compromiso, dos elementos hacen que esta sea más efectiva: (i) la credibilidad de la autoridad monetaria (ii) que el compromiso sea explicito y verificable ex post.
\end{la}

\clearpage
\section{Conclusión} 


\subsection{Recapitulación}


\subsection{Extensiones y relación con otras partes del Temario}



\subsection{Opinión}

\subsubsection{Idea final (Salida o cierra)}

\clearpage
\pagenumbering{gobble}
\MyBibliography



\MyBibItem{DAWID y GATTI}{Chapter 2 - Agent-Based Macroeconomics}{2018}{https://doi.org/10.1016/bs.hescom.2018.02.006}


% --- Anexos ---
\clearpage
\anexo{\textbf{Preguntas Test}}
%\input{\TemaCode/questions.tex} 


\clearpage
\anexo{Bloque Teoría Monetaria: Introducción}\label{anx:historia}
Con éste tema comienza el apartado del temario dedicado a Teoría monetaria. 

La introducción del dinero en macroeconomía puede entenderse como una extensión (imprescindible), ya que los modelos básicos asumen neutralidad nominal al partir del supuesto de precios perfectamente flexibles; como resultado, obvian la presencia del dinero en los sistemas económicos. Esto es así tanto en los modelos de crecimiento económico, como en los del ciclo real. En relación a los últimos, la evidencia empírica comprueba la existencia de rigideces nominales y reales, y motiva así tanto la introducción del dinero en los modelos, como la actuación de los bancos centrales para afectar a su cantidad en circulación y precio. 

El temario aborda la teoría monetaria de la siguiente manera: 
\begin{lnum}
    \item \authorfont{Tema 3.A.35}. El dinero puede entenderse no como un bien de consumo sino como un activo financiero que además genera una rentabilidad negativa (dado el coste de oportunidad de su tenencia). Por tanto, podemos preguntarnos ¿por qué demandarían dinero los agentes económicos? Y, de existir tal función de demanda, ¿cuáles son sus determinantes y con qué magnitud y condicionantes?
    \item \authorfont{Tema 3.A.37}. ¿Alteran los niveles (precio y cantidad, tipos y agregados, respectivamente) de dinero en una economía el equilibrio agregado? Es decir, ¿tiene el dinero efectos reales? ¿Cómo se transmiten los cambios en estos niveles al resto del sistema? ¿Qué dicen la teoría y la evidencia empírica?
    \item \authorfont{Tema 3.A.36}. Teniendo en cuenta los resultados de los dos anteriores temas, ¿cómo debe el banco central diseñar su política monetaria? Siendo conscientes de que el banco central no decide directamente la oferta de dinero, ¿cómo implementar esta política?
    \item \authorfont{Tema 3.A.39} y \authorfont{Tema 3.A.40}. El dinero y la deuda pública están estrechamente relacionados desde múltiples ópticas, ¿qué nos dice la teoría económica al respecto? ¿Y qué métodos empíricos existen para estudiar esta relación?
    \item \authorfont{Tema 3.A.41}. La política monetaria tiene efectos colaterales... en especial, la inflación;
    \item \authorfont{Tema 3.B.15} y \authorfont{Tema 3.B.16}. Hasta ahora, se ha estudiado la economía monetaria en un marco de economía cerrada, ¿qué ocurre cuando relajamos ese supuesto e introducimos el resto del mundo? 
    \item \authorfont{Tema 3.B.43}. ¿Cómo aterriza este cuerpo de conocimiento el BCE en el ejercicio de su mandato para el caso de la Eurozona?
\end{lnum}


\end{document}